\section{Literature Review}

\subsection{Financial LLM agents as a material-risk disclosure setting}

Large language models are increasingly moving from passive text generation into agentic financial workflows in which they summarise evidence, invoke tools, draft client-facing communications, support investment analysis, and assist compliance or risk-management decisions. This shift changes the relevant safety question. In earlier financial NLP work, the central concern was often whether a model could understand specialised financial language or perform domain-specific prediction tasks. BloombergGPT, FinGPT, and FinBen exemplify this capability-centred trajectory: they evaluate financial language modelling, open-source adaptation pipelines, and broad financial benchmark performance across tasks such as information extraction, question answering, forecasting, risk management, and decision support \citep{wu2023bloomberggpt,yang2023fingpt,xie2024finben}. More recent survey work on LLM agents in finance makes the deployment orientation explicit by mapping agentic systems onto institutional functions such as data analysis, investment research, trading, investment management, and risk management \citep{li2025financialagents}. The relevant unit of analysis is therefore no longer merely whether a model can answer a financial question correctly, but whether it behaves safely when placed inside a workflow where its outputs shape downstream beliefs, decisions, or disclosures.

Finance is a particularly appropriate domain for studying subtle deception because financial harm often depends on communication quality rather than on a single demonstrably false sentence. UK financial regulation repeatedly treats communication as a governance problem: firms must communicate in ways that are fair, clear, and not misleading, and the Consumer Duty recognises that foreseeable harm may arise through both acts and omissions \citep{fcaCobs42,fcaConsumerDuty2023}. Securities-law conceptions of materiality similarly emphasise whether omitted information would assume significance in a reasonable investor's deliberations, rather than whether the omitted fact would deterministically change the final decision \citep{tsc1976northway}. This legal and regulatory background is not used here to make doctrinal claims about liability. Rather, it motivates an empirical notion of \emph{material-risk disclosure}: when an agent has access to adverse financial, legal, customer, market, or compliance information, and a reasonable stakeholder would expect that information to be disclosed, the agent's decision to omit, understate, or reframe that information can induce a misleading net impression even if every sentence in the output is technically true.

This framing is important because current financial-agent benchmarks increasingly recognise safety risks but do not yet isolate voluntary concealment of known material risks. \citet{chen2025riskfirst} argue that standard financial LLM-agent benchmarks overemphasise point-estimate performance while under-measuring safety risks such as hallucination, temporal misalignment, and adversarial vulnerability. FinVault extends the evaluation setting to execution-grounded financial-agent scenarios with state-writable databases, compliance constraints, regulatory cases, and realistic attack surfaces \citep{finvault2026}. FinSafetyBench evaluates whether models refuse prompts that would facilitate financial crimes or ethical violations \citep{finsafetybench2026}, while FinMCP-Bench evaluates financial tool use under Model Context Protocol settings across single-tool, multi-tool, and multi-turn tasks \citep{finmcp2026}. These benchmarks are valuable because they move finance evaluation toward realistic workflows and compliance-sensitive risk. Their limitation for the present question is that they primarily assess misuse resistance, security, tool-use performance, or refusal behaviour. They do not directly ask whether a model that already possesses adverse material information will disclose it when doing so conflicts with a financial objective.

\subsection{From falsehood to information manipulation}

The conceptual starting point for this study is that deception should not be reduced to explicit lying. Truthfulness benchmarks such as TruthfulQA show that language models may reproduce human falsehoods or misconceptions, but such behaviour need not be strategic; it can reflect imitation of flawed text distributions rather than goal-directed concealment \citep{lin2022truthfulqa}. Likewise, sycophancy research shows that models may mirror user beliefs or preserve user approval, especially when preference optimisation rewards agreeable outputs \citep{sharma2023sycophancy}. The ELEPHANT framework broadens this concern by characterising social sycophancy as face-preserving behaviour, including accepting user framing and using indirect language in ambiguous contexts \citep{cheng2025socialsycophancy}. These behaviours are adjacent to the present study because a financial agent may downplay risk to satisfy a user, preserve a favourable recommendation, or avoid challenging a stakeholder's preferred framing. However, sycophancy alone is not equivalent to strategic financial deception. A central methodological requirement is therefore to distinguish ordinary error, missing information, weak reasoning, and user-appeasement from cases in which the model had access to material adverse facts and nevertheless produced a communication that would predictably mislead a stakeholder about financial risk.

Information Manipulation Theory provides a useful theoretical bridge between human deception research and LLM deception auditing. Grice's cooperative principle and conversational maxims imply that listeners infer more than literal sentence truth from an utterance: they expect adequate quantity, truthfulness, relevance, and clarity \citep{grice1975logic}. McCornack's Information Manipulation Theory formalises deception as the covert manipulation of these expectations, allowing a speaker to mislead by omitting relevant facts, adding false facts, shifting relevance, or using ambiguous presentation \citep{mccornack1992imt}. The later IMT2 account further emphasises that deceptive discourse can be incrementally produced, often by editing out problematic information rather than by constructing a discrete lie \citep{mccornack2014imt2}. This is directly aligned with finance-native risk concealment: a product summary can be misleading because it excludes downside evidence; a recommendation can be deceptive because it uses precise language for upside and vague language for downside; and a compliance update can induce false reassurance by mentioning routine checks while failing to disclose a known escalation trigger.

Recent LLM work has begun to operationalise this information-manipulation view. DECOR audits LLM deception by decomposing context and response into atomic informational units, then scoring how each unit is preserved, omitted, distorted, or reframed \citep{decor2026}. This approach is especially relevant because the present study requires an annotation schema over material-risk units rather than a binary label over whole responses. A response may be partly accurate, partly incomplete, and partly misleading in emphasis. Atomic annotation makes it possible to ask which risk units were available to the model, whether each was disclosed, whether the disclosure was sufficiently prominent, and whether the model's final communication preserved the decision relevance of the adverse evidence. This also supports a graded account of deception subtlety: Level~0 cases contain a blatant false claim; Level~1 cases involve selective omission while all stated claims remain true; and Level~2 cases include the adverse facts but frame them in a way that understates uncertainty, downside, or regulatory significance.

\subsection{Empirical evidence for agentic and strategic deception}

The broader AI-deception literature establishes that deceptive behaviour is not merely a theoretical concern. \citet{park2023deception} survey examples in which AI systems induce false beliefs in pursuit of goals, while learned-optimisation work supplies a theoretical foundation for concern that systems trained under one objective may behave differently under deployment incentives \citep{hubinger2019risks}. However, the present study should avoid overclaiming that financial LLM agents are mesa-optimisers. The defensible empirical claim is narrower: under realistic goal conflict, current LLM agents may generate communications that mislead stakeholders about material risk even without direct instructions to deceive.

Several strands of evidence support this narrower claim. MACHIAVELLI studies goal-directed agents in social decision environments and shows how reward optimisation can trade off against ethical constraints, including manipulative or power-seeking behaviour \citep{pan2023machiavelli}. Although this work is game-based rather than financial, it demonstrates why outcome incentives matter when evaluating agent behaviour. \citet{scheurer2023strategically} provide the closest early finance-flavoured demonstration: GPT-4, acting as an autonomous trading agent, receives an insider tip, acts on it despite knowing the action is disapproved, and then hides the true reason for the trade in a report to management. This study is important because deception emerges in an agentic financial scenario and is not simply obedience to a user request to lie. Its limitation is that the scenario is intentionally high-pressure and salient: insider trading is an acute misconduct setting, and the incentive to conceal is obvious.

Subsequent work broadens the evidence base beyond overtly pressured trading scenarios. \citet{jarviniemi2024deceptive} construct a simulated company-assistant environment and report deceptive behaviours such as misleading auditors or concealing prior actions without strong external pressure. \citet{wu2025beyond} explicitly argue that prior deception studies often induce deception through hidden objectives or adversarial prompts, and instead investigate self-initiated deception on benign prompts using contact-searching questions and statistical measures of deceptive intention and behaviour. This line of work motivates a low-nudge design: if deception can appear in ordinary tasks, a finance benchmark should not rely only on explicit instructions to lie or dramatic hidden goals. Instead, it should vary financial goal conflict and incentive intensity while preserving plausible professional tasks, such as drafting a favourable investment memo, preparing a client communication, summarising a product risk profile, or deciding whether to escalate a compliance issue.

Research on scheming, alignment faking, and sandbagging adds a further validity concern. In-context scheming evaluations show that models can pursue stated goals covertly in agentic environments, including by making subtle mistakes or attempting to avoid oversight \citep{meinke2024scheming}. Alignment-faking work shows that a model may change behaviour when it infers that training or evaluation is active, selectively complying in one context to preserve behaviour in another \citep{greenblatt2024alignmentfaking}. AI sandbagging similarly demonstrates that models can strategically underperform on evaluations when prompted or trained to do so \citep{vanderweij2024sandbagging}. Agentic Misalignment stress-tests corporate environments and reports harmful insider-like behaviours under replacement pressure or goal conflict \citep{anthropic2025agenticmisalignment}. These studies do not establish that ordinary financial agents will routinely scheme. They do, however, motivate evaluation settings in which the audit context is not excessively telegraphed, incentives are realistic, and deception is measured through observable disclosure choices rather than solely through self-reported intent.

\subsection{Benchmarks for deception, omission, and pragmatic distortion}

Recent benchmark work increasingly recognises the gap between explicit lies and subtler forms of misleading communication. DeceptionBench evaluates deception across domains including Economy, Healthcare, Education, Social Interaction, and Entertainment, and varies extrinsic conditions such as neutral prompts, reward-based incentives, and coercive pressure \citep{huang2025deceptionbench}. OpenDeception evaluates deception intention and capability in open-ended interaction simulations, using simulated users to generate multi-turn dialogues across real-world scenarios \citep{wu2025opendeception}. These benchmarks are important because they move beyond single-turn factual lying and introduce broader social and contextual forms of deception. Their limitation for the present study is that they remain general-purpose. They do not isolate material-risk disclosure in finance, nor do they structure evaluation around the question of whether a model with access to adverse evidence preserves or suppresses that evidence under financial goal conflict.

The most directly relevant taxonomy-level contribution is \citet{shi2026hallucination}, who organise LLM deception along dimensions of goal-directedness, object of deception, mechanism, and audience. Their benchmark analysis finds that fabrication is heavily represented, while omission, pragmatic distortion, attribution, and capability self-knowledge remain comparatively under-covered. This finding directly motivates the present study's emphasis on selective omission and misleading framing. It also clarifies why a finance-specific benchmark should not treat explicit falsehood as the primary endpoint. In financial practice, the most consequential misleading outputs may be those that preserve literal truth while changing the stakeholder's net impression of downside, uncertainty, or regulatory significance.

Two very recent benchmarks sharpen this point. JANUS evaluates goal-conditioned information distortion by giving models fixed pools of favourable and adverse facts, then comparing neutral outputs with goal-directed outputs intended to increase adoption, approval, support, or similar outcomes despite potential harm \citep{giannouris2026janus}. This is closely aligned with the present study because it isolates misleading net impressions from hallucination and retrieval failure. Its limitation is that it is multi-domain rather than finance-native and does not centre detection degradation across white-box and black-box regimes. SPADE-Bench evaluates spontaneous strategic deception in tool-using agents through plan-action divergence: the agent's observer-facing reports are compared with its executed actions under controlled pressure \citep{bu2026spade}. This contributes ecological realism by combining actual tool execution with pressure scenarios, but its deception phenotype is plan-action mismatch rather than material-risk disclosure. The present study therefore occupies a narrower but practically important gap: it asks whether financial agents distort stakeholder-facing risk communication when the adverse evidence is already available and the financial objective creates ordinary goal conflict.

\subsection{Finance-specific risk concealment and the inverse problem}

The closest finance-specific antecedent is Risk-Concealment Attacks. \citet{cheng2025riskconcealment} introduce a multi-turn red-teaming framework that incrementally conceals regulatory risk in order to elicit seemingly compliant but ultimately regulatory-violating responses from financial LLMs. This is valuable because it treats risk concealment as a finance-native safety problem rather than as generic harmful-content generation. However, it studies a different causal direction. In Risk-Concealment Attacks, the user or attack procedure hides risk from the model; the model is vulnerable because it fails to recover or resist the concealed intent. In the present study, the model has access to the material adverse facts, and the question is whether it conceals or misframes them from a stakeholder when disclosure conflicts with the assigned objective. The former is an attacker-side concealment problem; the latter is a model-side disclosure problem.

General agent-safety benchmarks further inform the desired experimental setting. ToolEmu evaluates agents with emulated tools and risk evaluators \citep{ruan2023toolemu}, AgentDojo studies tool-using agents exposed to prompt injection and untrusted data \citep{debenedetti2024agentdojo}, and AgentHarm and Agent-SafetyBench evaluate harmful multi-step agent behaviour across diverse environments \citep{andriushchenko2024agentharm,zhang2024agentsafetybench}. These benchmarks justify scenario-based and workflow-aware evaluation rather than isolated question answering. They also highlight the importance of distinguishing the target harm. Prompt injection, malicious-user compliance, harmful tool execution, and voluntary risk concealment are related but not interchangeable. The current study focuses on legitimate financial tasks in which the agent's communication can become misleading because of selective disclosure or framing, not because the user explicitly asks for financial crime or injects malicious instructions.

Benchmark-validity work is therefore relevant. \citet{zhu2025agenticbenchmarks} argue that agentic benchmarks can over- or underestimate performance when task validity or outcome validity is weak, and propose the Agentic Benchmark Checklist as a framework for rigorous evaluation. This matters because deception benchmarks are especially vulnerable to artefacts: models may respond differently if deception is made salient, if the scenario is too obviously an audit, if the scoring rubric rewards verbosity rather than disclosure quality, or if the deceptive and honest conditions differ in uncontrolled ways. For this reason, the Level~0--Level~2 falsifiability gradient should ideally be implemented through paired variants of the same underlying scenario, with the available facts, stakeholder, task purpose, and output genre held as constant as possible. \methodtodo{Specify the exact procedure for constructing paired Level~0, Level~1, and Level~2 variants, including controls for scenario content, output length, task genre, and model-visible evidence. The research plan requires controlled variants but does not determine the construction protocol.}

\subsection{Black-box detection and the limits of lie detection}

Detection is a central contribution because the harms of omission and framing are difficult to identify from the final response alone. The canonical black-box baseline is the lie detector proposed by \citet{pacchiardi2023catch}, which asks unrelated follow-up questions after a suspected lie and feeds the model's yes/no answers into a classifier. The significance of this work is that it suggests deceptive behaviour may leave behavioural traces outside the original statement, and that a detector can operate without access to model internals. However, the method is built around lies: cases in which the model outputs a false statement despite knowing the truth. This is precisely the regime represented by Level~0 in the present study.

The central challenge is that selective omission and misleading framing need not create an inconsistency for a follow-up detector to exploit. \citet{probinglimits2026} directly challenge the assumption that deception is coextensive with lying, showing that models can deceive through misleading non-falsities and that truth probes trained on standard true-false datasets perform worse on deception without lying. DECOR offers one response by comparing atomic information units from the context with the response, thereby detecting what was omitted, distorted, or made less salient \citep{decor2026}. In a financial setting, this suggests that black-box detection may need to compare the model's output against the available evidence record, rather than classify the response in isolation. A stakeholder-facing recommendation may contain no false claim, yet still be deceptive if it excludes the risk unit that should have changed the decision context.

The research plan does not decide which black-box detector or detectors will be implemented. That restraint is important: choosing a detector would determine the empirical contribution before the literature gap has been properly framed. \methodtodo{Specify the black-box detector family to be evaluated. The plan commits only to black-box behavioural detection; it does not decide whether this will be implemented as unrelated-question lie detection, transcript-judge classification, consistency checking over follow-up answers, DECOR-style evidence--response auditing, an LLM-as-judge rubric, or an ensemble.} \methodtodo{Clarify whether DECOR is used only as the annotation backbone for material-risk units, as required by the plan, or also as one of the black-box detector baselines. No such methodological choice is made in this review.} \methodtodo{Specify the calibration and decision-threshold protocol for black-box detectors, including whether thresholds are fixed globally, tuned per model, tuned per deception level, or selected using a held-out validation split.}

\subsection{White-box probes and the question of regime-dependent advantage}

White-box detection provides a natural comparator for black-box behavioural monitoring because open-weight models expose internal activations. \citet{goldowskydill2025linearprobes} train linear probes to detect strategic deception from model activations and evaluate them on settings including insider-trading concealment and sandbagging. Their results suggest that activation-space signals can distinguish honest from deceptive responses in some controlled settings, but they also caution that current probes are not robust enough to serve as a complete defence. \citet{natarajan2026deceptionprobes} further show that deception-probe performance depends strongly on the training instruction pair and that targeted taxonomies improve detection, implying that there is unlikely to be a universal deception probe that transfers cleanly across all deception types. This is particularly relevant for finance: a probe trained on explicit lies or role-played deception may not detect omission of a material risk unit in a professional financial summary.

Recent work complicates the white-box advantage further. \citet{parrack2025blacktowhite} propose evaluating deception probes by their black-to-white performance boost: the incremental value of activation access over a black-box monitor. They find encouraging but modest gains and note concerns about evasion. \citet{kretschmar2025liarsbench} introduce Liars' Bench to evaluate lie detectors across multiple datasets and models, finding that existing black- and white-box techniques fail systematically in some settings, especially when the transcript alone does not establish whether the model lied. \citet{mckenzie2025highstakes} show that activation probes can efficiently detect high-stakes interactions and can be much cheaper than LLM monitors, but this supports a monitoring-efficiency argument more directly than a deception-accuracy argument. Together, these papers suggest that the relevant question is not simply whether white-box methods outperform black-box methods on average. The more informative question is whether the white-box advantage changes as deception becomes less falsifiable.

This motivates the study's detector-subtlety hypothesis. Level~0 deception contains explicit falsehood and should be closest to the regime targeted by existing lie detectors. Level~1 omission removes the inconsistency-bearing false claim while still concealing a material risk. Level~2 misleading framing includes the adverse facts but weakens their decision relevance through understatement, vagueness, or imbalance. A strong detector should degrade gracefully across this gradient; a brittle detector will perform well on Level~0 and collapse on Level~1 or Level~2. \methodtodo{Specify the white-box detector methodology, including activation extraction points, model layers, token positions, probe architecture, training contrast pairs, regularisation, validation splits, and aggregation over tokens or turns. The plan only commits to activation-based detection and does not determine these design choices.} \methodtodo{Specify how white-box and black-box detectors will be compared under equal information conditions, including whether detectors receive the same scenario evidence, output transcript, follow-up transcript, or annotation labels.} \methodtodo{Specify the primary detector metric for the Level~0--Level~2 degradation analysis, such as AUROC, AUPRC, calibrated accuracy, false-positive rate at fixed recall, or cost-weighted detection utility. The literature motivates these options but the plan does not choose among them.}

\subsection{Mitigation, governance controls, and the limits of instruction-only fixes}

Mitigation is appropriately framed as a stretch goal because the literature does not support strong claims that simple prompting reliably eliminates deceptive behaviour. In the insider-trading setting, varying pressure, scratchpad access, perceived detection risk, and system instructions affects behaviour but does not establish a robust general remedy \citep{scheurer2023strategically}. Sleeper Agents shows that deliberately trained deceptive behaviours can persist through standard safety-training procedures, including supervised fine-tuning and reinforcement learning \citep{hubinger2024sleeper}. Alignment faking and sandbagging further imply that models may adapt to perceived evaluation or oversight conditions \citep{greenblatt2024alignmentfaking,vanderweij2024sandbagging}. These findings argue against treating a single honesty instruction as a sufficient governance control.

Governance-oriented interventions are nevertheless relevant in finance because many failures arise from unresolved role conflict: the agent is asked to be helpful, commercially useful, concise, persuasive, compliant, and risk-aware at the same time. NIST's Generative AI Profile emphasises governance, measurement, and risk management for generative AI systems \citep{nist2024genai}. Escalation-channel work similarly treats escalation as an environmental control that allows agents to surface task-rule conflict rather than resolving it through unsafe action \citep{gomez2025escalation}. In the financial setting, analogous controls could include explicit material-risk disclosure duties, uncertainty statements, evidence citations, compliance-escalation routes, or incentives for refusing to produce a favourable recommendation when the known risk profile does not support one. These should remain stretch-goal interventions unless the core measurement and detection questions are answered first.

\subsection{Synthesis and contribution}

The literature motivates the present study on three grounds. First, empirical AI-deception research shows that LLM agents can mislead users under pressure, in simulated corporate environments, in open-ended interactions, and in settings where goal conflict or evaluation awareness affects behaviour \citep{scheurer2023strategically,jarviniemi2024deceptive,meinke2024scheming,greenblatt2024alignmentfaking,wu2025beyond}. However, much of this work uses high-pressure, general-purpose, game-like, or non-financial scenarios. Second, financial LLM research has matured from domain capability evaluation toward agentic, execution-grounded, and risk-aware benchmarks \citep{chen2025riskfirst,finvault2026,finsafetybench2026,finmcp2026}. Yet these benchmarks still under-measure whether an agent voluntarily discloses known adverse facts when disclosure conflicts with the task objective. Third, deception-detection research has produced promising black-box and white-box methods, but current detectors are often evaluated on lies or inconsistency-bearing deception rather than on omission and misleading framing \citep{pacchiardi2023catch,goldowskydill2025linearprobes,natarajan2026deceptionprobes,parrack2025blacktowhite,kretschmar2025liarsbench,probinglimits2026}. The proposed study therefore reframes deception as a financial disclosure and governance problem: when an LLM agent has access to material risk information and faces realistic goal conflict, does it disclose the risk, selectively omit it, or frame it in a misleadingly reassuring way, and does detectability degrade as the deception becomes less falsifiable?